% !TeX root = ../../Book3.tex
% 纲要标题：### 21.6* 分片多项式与样条计算
% * 直线两侧的光滑拼接与差多项式被直线方程的幂整除
% * 拼接矩阵、辅助变量和合冲模；明确工作环及次数滤过
% * 四象限实例中全部解的自由模表示、顶点核验及次数受限的维数
% * 只保留这一小型代数应用；整数规划、多面体与一般网格理论暂不扩入本卷
% ---

\OptionalSection{分片多项式与样条计算}{b3:sec:2106}

给定一块分成若干区域的平面，分别在各区域选取多项式，还要使它们在公共边上连续或具有相同的若干阶导数。区域内的系数可以自由选择，拼接条件却把不同区域联系起来。下面把这些条件写成多项式矩阵，并用合冲模求出全部允许的选择。本节只讨论平面上直线分隔的简单区域，不涉及逼近误差或一般网格的维数公式。

\subsection{光滑拼接与整除条件}
\label{b3:subsec:2106:01}

\begin{definition}\label{b3:def:polynomial-spline}
设平面区域 \(\Omega\subseteq\mathbb R^2\) 由有限个闭多边形区域 \(\Delta_1,\ldots,\Delta_s\) 拼成；允许这些区域无界，其边可以是线段、射线或直线。各片内部互不相交，交集为公共边、顶点或空集。给定整数 \(r\ge0\)，若函数 \(p:\Omega\to\mathbb R\) 在每片 \(\Delta_i\) 上由多项式 \(p_i\in\mathbb R[x,y]\) 给出，且这些分片在 \(\Omega\) 的内部拼成 \(C^r\) 函数，则称 \(p\) 为该划分 \(\Delta=(\Delta_1,\ldots,\Delta_s)\) 上的 \(C^r\) \term[duoxiangshi-yangtiao]{多项式样条}[polynomial spline]。对整数 \(d\ge0\)，所有分片次数不超过 \(d\) 的样条组成的实向量空间记为 \(C_d^r(\Delta)\)。
\end{definition}

先只考虑一条公共边的内部。写它所在直线为 \(\ell=ax+by+c=0\)，其中 \((a,b)\ne(0,0)\)。沿边的切向坐标和横向坐标 \(\ell\) 构成一个可逆仿射坐标变换。

\begin{proposition}\label{b3:prop:spline-divisibility}
设 \(p_+,p_-\in\mathbb R[x,y]\)，\(\ell=ax+by+c\in\mathbb R[x,y]\)，其中 \((a,b)\ne(0,0)\)，并取整数 \(r\ge0\)。若两个平面区域在直线 \(\ell=0\) 的一段非空开线段两侧相接，则由 \(p_+\)、\(p_-\) 给出的分片函数在该线段附近拼成 \(C^r\) 函数，当且仅当
\[
p_+-p_-\in(\ell^{r+1}).
\]
\end{proposition}
\begin{proof}
取仿射坐标 \((u,v)\)，使 \(v=\ell\)，并写
\(p_+-p_-=\sum_{j=0}^N a_j(u)v^j\)。若能拼成 \(C^r\)，则沿线段的前 \(r\) 阶横向导数相等，故 \(j!a_j(u)=0\) 在一个开区间上成立，\(0\le j\le r\)。实多项式在区间上恒零必为零多项式，且 \(j!\ne0\)，所以前 \(r+1\) 项全部为零，得到整除性。

反之，若差为 \(v^{r+1}q(u,v)\)，取总阶数不超过 \(r\) 的任意偏导数，所得每项仍含 \(v\) 因子，所以两侧的这些导数在线段上相等。先由函数值相等拼接连续函数，再逐阶利用一元差商说明横向导数在界面存在，切向导数由界面上的相同多项式给出。归纳得到 \(C^r\) 拼接。
\end{proof}

例如沿 \(x=0\) 作 \(C^1\) 拼接，全部两片式样条可唯一写成
\[
(p_-,p_+)=(p,p+x^2q),\qquad p,q\in\mathbb R[x,y].
\]
因此给出一片多项式 \(p\) 和一个任意多项式 \(q\)，就给出了全部可能的拼接；只令两片在原点取值相等显然不足以得到这个结论。

\subsection{拼接条件的矩阵表示}
\label{b3:subsec:2106:02}

令 \(S=\mathbb R[x,y]\)。为每条内部公共边选择一个方向 \(i\to j\)，并以 \(\ell_{ij}=0\) 表示其直线。引入辅助多项式 \(q_{ij}\)，把边上的条件写成
\[
p_j-p_i-\ell_{ij}^{r+1}q_{ij}=0.
\]
若共有 \(e\) 条内部边，这些等式组成一个 \(e\times(s+e)\) 多项式矩阵 \(M\)。边条件解模就是 \(\ker M\) 在前 \(s\) 个坐标上的投影。因 \(S\) 为整环且 \(\ell_{ij}\ne0\)，给定 \((p_i)\) 后各 \(q_{ij}\) 唯一，所以这个投影在核上单射。

\ProseParagraphBreak
环 \(S\) 的 Noether 性保证 \(\ker M\) 有有限生成组。实际调用\cref{b3:alg:polynomial-matrix-kernel}时，还须让边方程的系数属于一个支持精确四则运算与零判定的实数子域 \(k\)，例如 \(\mathbb Q\) 或给定的实代数数域。先在 \(k[x,y]\) 上求核，再扩张系数到 \(\mathbb R\)，便得到全部实系数关系：域扩张 \(\mathbb R/k\) 为平坦扩张，因而核的正合列在扩张系数后仍然正合。有限生成组的投影于是给出所有满足边条件的分片多项式。任意实数若只以近似值输入，并不自动满足这个算法条件。

\ProseParagraphBreak
对只有一条分隔线的区域，\cref{b3:prop:spline-divisibility}已经证明这些条件就是样条条件。若若干边在内部顶点相遇，还须核对顶点处的拼接；下节的具体区域将逐项处理这一点，而不从边的内部条件直接跳到任意网格的一般定理。

\ProseParagraphBreak
次数限制也需要单独记录。仿射直线含常数项时，解模通常只有次数滤过；通过原点的直线则使关系齐次，可以把辅助坐标赋予次数平移，从而使用本章的分次 Hilbert 计算。仅数矩阵的行数与列数，不能求出给定次数下的维数，因为不同边的条件可能存在关系。

\subsection{四个象限上的完整计算}
\label{b3:subsec:2106:03}

用坐标轴把平面分为四个闭象限，依次编号为右上、左上、左下、右下。记分片为 \((p_1,p_2,p_3,p_4)\)，并令 \(s=r+1\)。沿四条半轴的条件写为
\[
\begin{aligned}
p_2-p_1&=x^s a,& p_3-p_2&=y^s b,\\
p_4-p_3&=x^s c,& p_1-p_4&=y^s d.
\end{aligned}
\]
相加得到唯一的环行相容条件
\[
x^s(a+c)+y^s(b+d)=0.
\]
因为 \(x^s,y^s\) 互素，全部关系为
\[
a+c=y^s h,\qquad b+d=-x^s h\qquad(h\in S).
\]
证明也很直接：\(x^s\mid y^s(b+d)\) 推出 \(x^s\mid b+d\)，写 \(b+d=-x^sh\) 后再约去 \(x^s\)。这正是两生成元 \((x^s,y^s)\) 的合冲计算。

\ProseParagraphBreak
取 \(p=p_1\)，依次回代得到全部解
\[
\begin{aligned}
(p_1,p_2,p_3,p_4)
={}&p(1,1,1,1)+a(0,x^s,x^s,0)\\
&+b(0,0,y^s,y^s)+h(0,0,0,x^sy^s).
\end{aligned}
\]
这四个向量的系数唯一：先读第一坐标求 \(p\)，再由第二、第三、第四坐标顺次求 \(a,b,h\)。所以边条件解模是自由分次模
\[
S\oplus S(-s)^2\oplus S(-2s).
\]
四个生成元确实在原点也是 \(C^r\) 的。第二个对应左半平面上的 \(x^s\)、右半平面上的零，第三个对应下半平面上的 \(y^s\)、上半平面上的零；第四个是右半平面的分片 \(x^s\) 与下半平面的分片 \(y^s\) 的乘积。每个一元分片函数的前 \(r=s-1\) 阶导数都在零处连续，乘以全局多项式及有限相加保持 \(C^r\)。因此上述全部解恰为整个平面的样条，顶点处没有遗漏的条件。

\ProseParagraphBreak
记 \(N(j)=\binom{j+2}{2}\) 当 \(j\ge0\)，而 \(N(j)=0\) 当 \(j<0\)。它是二元全次数不超过 \(j\) 的多项式空间维数。上述唯一表示逐项保持次数，故
\[
\dim_{\mathbb R}C_d^r(\Delta)=N(d)+2N(d-s)+N(d-2s).
\]
例如连续分片二次多项式对应 \(r=0,d=2\)，维数为 \(6+2\cdot3+1=13\)；若要求一阶光滑，则 \(r=1\)，维数变为 \(6+2=8\)。环行条件中的辅助元 \(h\) 正是四条边约束之间的关系所留下的自由量。这说明为什么实际维数不能只用四片的系数总数减去表面上列出的方程数。
